\documentclass[12pt,oneside]{book}

%============================================================
% PACKAGES
%============================================================

%--------------------
% Page + font
%--------------------
\usepackage[margin=1in]{geometry}
\usepackage[T1]{fontenc}
\usepackage[utf8]{inputenc}
\usepackage{lmodern}
\usepackage{microtype}
\usepackage{parskip}


%--------------------
% Math
%--------------------
\usepackage{amsmath, amssymb, amsfonts, amsthm, mathtools}
\usepackage{bm}
\usepackage{bbm}
\usepackage{mathrsfs}
\usepackage{dsfont}

%--------------------
% Tables + figures
%--------------------
\usepackage{graphicx}
\usepackage{float}
\usepackage{booktabs}
\usepackage{array}
\usepackage{xltabular}
\usepackage{multirow}
\usepackage{caption}
\usepackage{subcaption}

%--------------------
% Lists
%--------------------
\usepackage{enumitem}
\setlist[itemize]{topsep=2pt,itemsep=2pt}
\setlist[enumerate]{topsep=2pt,itemsep=2pt}

%--------------------
% Colors + hyperlinks
%--------------------
\usepackage[dvipsnames]{xcolor}
\usepackage[
    colorlinks=true,
    linkcolor=MidnightBlue,
    citecolor=MidnightBlue,
    urlcolor=MidnightBlue
]{hyperref}
\usepackage[nameinlink,capitalise]{cleveref}

%--------------------
% Algorithms
%--------------------
\usepackage{algorithm}
\usepackage{algpseudocode}

%--------------------
% Code blocks
%--------------------
\usepackage{listings}

\lstdefinestyle{codestyle}{
    basicstyle=\ttfamily\small,
    breaklines=true,
    frame=single,
    numbers=left,
    numberstyle=\tiny,
    keywordstyle=\color{MidnightBlue},
    commentstyle=\color{ForestGreen},
    stringstyle=\color{BrickRed},
    showstringspaces=false,
    tabsize=4
}

\lstset{style=codestyle}

%--------------------
% Section styling
%--------------------
\usepackage{titlesec}
\usepackage{tocloft}

\setcounter{tocdepth}{2}
\setcounter{secnumdepth}{3}

%============================================================
% THEOREMS
%============================================================

\numberwithin{equation}{section}

\theoremstyle{definition}
\newtheorem{definition}{Definition}[section]
\newtheorem{example}{Example}[section]
\newtheorem{exercise}{Exercise}[section]
\newtheorem{remark}{Remark}[section]

\theoremstyle{plain}
\newtheorem{theorem}{Theorem}[section]
\newtheorem{lemma}{Lemma}[section]
\newtheorem{proposition}{Proposition}[section]
\newtheorem{corollary}{Corollary}[section]

\theoremstyle{remark}
\newtheorem*{note}{Note}

%============================================================
% CUSTOM COMMANDS
%============================================================

% Sets
\newcommand{\R}{\mathbb{R}}
\newcommand{\N}{\mathbb{N}}
\newcommand{\Z}{\mathbb{Z}}
\newcommand{\Q}{\mathbb{Q}}
\newcommand{\C}{\mathbb{C}}
\newcommand{\powerset}{\mathcal{P}}


% Probability
\newcommand{\E}{\mathbb{E}}
\newcommand{\Var}{\mathrm{Var}}
\newcommand{\Cov}{\mathrm{Cov}}
\newcommand{\Corr}{\mathrm{Corr}}
\newcommand{\Prob}{\mathbb{Pr}}

\newcommand{\ind}{\mathbbm{1}}
\newcommand{\iid}{\overset{\text{iid}}{\sim}}
\newcommand{\convd}{\overset{d}{\to}}
\newcommand{\convp}{\overset{p}{\to}}
\newcommand{\as}{\overset{a.s.}{\to}}

% Operators
\newcommand{\given}{\,\middle|\,}
\newcommand{\argmin}{\operatorname*{arg\,min}}
\newcommand{\argmax}{\operatorname*{arg\,max}}

% Linear algebra / analysis
\newcommand{\norm}[1]{\left\lVert #1 \right\rVert}
\newcommand{\abs}[1]{\left\lvert #1 \right\rvert}
\newcommand{\inner}[2]{\left\langle #1, #2 \right\rangle}

\newcommand{\dd}{\,\mathrm{d}}
\newcommand{\grad}{\nabla}
\newcommand{\Hess}{\nabla^2}

\newcommand{\tr}{\operatorname{tr}}
\newcommand{\rank}{\operatorname{rank}}
\newcommand{\diag}{\operatorname{diag}}

% Distributions
\newcommand{\Normal}{\mathcal{N}}
\newcommand{\Bern}{\mathrm{Bernoulli}}
\newcommand{\Bin}{\mathrm{Binomial}}
\newcommand{\Pois}{\mathrm{Poisson}}
\newcommand{\Gam}{\mathrm{Gamma}}
\newcommand{\BetaDist}{\mathrm{Beta}}
\newcommand{\Exp}{\mathrm{Exponential}}
\newcommand{\Dir}{\mathrm{Dirichlet}}
\newcommand{\InvGam}{\mathrm{Inv\text{-}Gamma}}

%============================================================
% TITLE INFO
%============================================================

\title{EN.553.740 Machine Learning 1}
\author{Jonathan Ma}
\date{\today}

%============================================================
% DOCUMENT
%============================================================

\begin{document}

\maketitle
\tableofcontents

\section{Introduction}
A good way to think of machine learning is instead of a computer taking data and a program to produce output, it takes data and output and attempts to produce a program, ie "learn" a function.

We use it for tasks that are too complex to program, like driving or large scale data analysis. A bit of history of "learning":
\begin{itemize}
    \item 1959: Arthur Samuel (IBM) created a program to play checkers, introducing rule-based learning.
    \item 1960: Frank Rosenblatt introduced a linear classifier based on brain neurons, called a Perceptron.
    \item 1969: Marvin Minsky showed limitations of the perceptron, halting development for 20 years
    \item 1980: First workshop in machine learning
    \item 1986: Multilayer networks and backpropagation introduced
    \item 1990: Vapnik and Chervonenkis introduce kernels and SVM, and termed "Statistical Learning"
    \item 2002: Rise of graphical models
    \item 2010: Deep learning takes over
    \item 2016: ML hits record size with NIPS and ICML doubling in size
\end{itemize}

Machine learning has been used for a variety of tasks, like stock market prediction, decoding thoughts from brain scans, and self-driving cars.

\subsection{Learning Problems}
Statistical Learning falls under five rough fields:
\begin{itemize}
    \item Classification (Assign a category)
    \item Regression (predict a real value)
    \item Ranking (order items)
    \item Clustering (item partitioning, like image segmentation)
    \item Dimensionality Reduction/Manifold Learning (Computer vision)
\end{itemize}
In a broader sense, there is supervised learning (Classification/regression) and unsupervised learning (Clustering, dimensionality reduction, density estimation) and things like semi-supervised learning, active learning, reinforcement learning, etc.

\paragraph{Supervised Learning}
For a given problem, we learn a function $f: X \rightarrow Y$ such that $x \in \mathcal{X}, \quad y \in \mathcal{Y}$. Classification deals with discrete labels while regression deals with continuous functions.

\paragraph{Unsupervised Learning}
For a given problem, we learn a function $f: X \rightarrow Y$ such that $x \in \mathcal{X}, \quad f(x) \in \R$. Usually embedding (for text), clustering, or density estimation.

\section{Mathematical Review}
Algorithms are built from sets, functions, vectors, matrices, norms, and derivatives, so we introduce common notation here. Consider a student success dataset with the target being quiz score, and features as study time and prior score.

\subsection{Sets and Power Sets}
Let $A = \{1,2,3,4,5\}$ be the set of students in the toy dataset. The \emph{power set} is $\powerset(A)=\{B: B \subseteq A\}$. If $|A|<\infty$, then $|\powerset(A)|=2^{|A|}$. A subset $B \subseteq A$ may represent a mini-batch, validation set, or a selected group of samples.

\subsection{Functions}
For sets $A,B$, write $B^A=\{f: A\rightarrow B\}$. if $B=\R$, then $\R^A$ is the set of real-valued functions on $A$. If $|A| < \infty$, then $\R^A \simeq \R^{|A|}$ as vector spaces. For our toy example: \[x_1: A \rightarrow \R \qquad x_2: A \rightarrow \R \qquad y: A \rightarrow \R\]
For a subset $C \subseteq A$, define an \emph{indicator function}:
\[
\mathbbm{1}_C: A \rightarrow \{0,1\}, \quad \mathbbm{1}_C(i)=
\begin{cases}
    1, i \in C, \\
    0, i \notin C
\end{cases}
\]
Indicators are used to encode events, groups, labels, masks, etc.

\subsection{Vectors}

\end{document}